\documentclass[conference]{IEEEtran}

\usepackage{graphicx}
\usepackage{amsmath}
\usepackage{amssymb}
\usepackage{booktabs}
\usepackage{multirow}
\usepackage{url}
\usepackage{cite}

% ─── IMAGES TO ADD (place these files in the same folder as this .tex) ────────
%
%  REQUIRED (paper will not compile without these):
%  1. system_architecture_paper.png  — block diagram:
%     [Images] → [Preprocessing] → [ResNet50] → [Features] → [KNN] → [Results]
%
%  2. Precision Recall F1.png — bar chart comparing Accuracy/F1 of 3 models
%
%  OPTIONAL (include if available from your Streamlit app screenshots):
%  3. Master Category Distribution.png  — from Analytics Tab 1
%  4. Confusion Matrix.png              — from Analytics Tab 2
%  5. Free Space-PCA Projection.png     — from Analytics Tab 1
%  6. tSNE_Visualization.png            — from models/tsne_visualization.png
%  7. Top-K Similarity Distributions.png — from Analytics Tab 3
%  8. Streamlit interface.png           — screenshot of Browse Catalog mode
%  9. Streamlit interface1.png          — screenshot of Upload & Discover mode
% ─────────────────────────────────────────────────────────────────────────────

\title{Visual Fashion Recommendation Using Deep Feature Embeddings:\\
A Three-Model Comparative Study with Interactive Analytics}

\author{
\IEEEauthorblockN{K. Teja,
                  A. Shwetha,
                  V. Shirisha,
                  Ch. Amrutha,
                  Dr. Shivadi Balakrishna}
\IEEEauthorblockA{Department of Computer Science
and Engineering (Advanced Computing)\\
Vignan's Foundation for Science, Technology and Research,
Andhra Pradesh, India\\
\{teja, shwetha, shirisha, amrutha, balakrishna.shivadi\}@vignan.ac.in}
}

\begin{document}
\maketitle

% ─────────────────────────────────────────────────────────────────────────────
\begin{abstract}
Keyword-based search fails to capture the visual nuances that drive
fashion purchases --- a shopper looking for ``something like this blue
kurta'' needs a system that reasons about shape, color, and texture,
not just product tags. We tackle this gap by building a content-based
visual recommendation pipeline on the Myntra catalog (44,441 product
images across seven categories) and running a head-to-head comparison
of three frozen ImageNet backbones under identical preprocessing,
indexing, and evaluation conditions. ResNet18 yields 512-dimensional
embeddings, VGG16 yields 4096-dimensional embeddings, and ResNet50
yields 2048-dimensional embeddings; all vectors are L2-normalized so
that a standard Ball Tree index can perform cosine-equivalent retrieval
without a GPU at query time. To prevent the dominant Apparel class
(48\% of the catalog) from inflating aggregate scores, we adopt a
stratified evaluation protocol that draws queries equally from each of
the seven master categories. On this balanced benchmark, ResNet50
achieves the strongest retrieval quality (Accuracy 0.777, F1-Score@6
0.789, Top-1 Cosine Similarity 0.940), while VGG16 produces the most
diverse recommendation lists (Diversity 0.163) at the cost of lower
precision. We wrap the full pipeline in a Streamlit web application
that supports catalog browsing, image-upload search, and a three-tab
analytics dashboard, demonstrating that the system is deployable with
modest computational resources.
\end{abstract}

\begin{IEEEkeywords}
fashion recommendation, content-based retrieval, ResNet50, VGG16,
ResNet18, transfer learning, KNN, Myntra dataset, Streamlit
\end{IEEEkeywords}


% ─────────────────────────────────────────────────────────────────────────────
\section{Introduction}

Indian fashion e-commerce platforms such as Myntra and Flipkart host
catalogs spanning tens of thousands of products --- from ethnic kurtas
to western sneakers, from personal care items to sporting goods. When a
shopper spots a colleague wearing a particular style of printed kurti,
typing ``blue printed kurti with mirror work'' into a search bar rarely
surfaces the right results. The gap between visual intent and textual
query is exactly where content-based image retrieval becomes useful
\cite{mcauley2015image}.

Traditional recommendation approaches rely on collaborative filtering
or purchase-history signals, both of which ignore the visual dimension
of fashion. Color palettes, silhouettes, fabric textures, and decorative
patterns are the features that actually guide a purchase decision, yet
none of these are easy to express as keywords. Deep convolutional
networks trained on ImageNet have been shown to learn visual
representations that transfer surprisingly well to downstream tasks,
including fashion retrieval, without domain-specific fine-tuning
\cite{yosinski2014transferable}.

In this work we build a complete retrieval pipeline on the publicly
available Myntra fashion dataset and conduct a controlled experiment
comparing three backbone architectures. Concretely, we make the
following contributions to the visual fashion recommendation
literature:

\begin{itemize}
  \item A controlled three-way comparison of ResNet18, VGG16, and
        ResNet50 under identical preprocessing, indexing, and evaluation
        conditions on 44,441 fashion images, reporting classification
        metrics, similarity scores, diversity, and catalog coverage.
  \item A stratified evaluation protocol that draws an equal number of
        queries from each of seven categories, ensuring that
        minority-class performance (e.g., Sporting Goods with only 25
        items) is not hidden behind the dominant Apparel class.
  \item An interactive Streamlit application covering catalog browsing,
        real-time image-upload search, and a three-tab analytics
        dashboard that visualizes dataset statistics, classification
        metrics, and recommendation quality indicators.
\end{itemize}

The remainder of this paper is organized as follows.
Section~\ref{sec:related} reviews related work. Section~\ref{sec:method}
describes the methodology. Section~\ref{sec:setup} covers the dataset
and evaluation setup. Section~\ref{sec:results} presents results and
discussion. Section~\ref{sec:app} describes the web application.
Section~\ref{sec:conclusion} concludes with limitations and future
directions.


% ─────────────────────────────────────────────────────────────────────────────
\section{Related Work}
\label{sec:related}

The trajectory from hand-crafted descriptors to deep features in visual
retrieval follows a clear arc. Early systems relied on SIFT
\cite{lowe2004distinctive} and HOG for local pattern matching, but
these features struggle with the deformable geometry of draped fabric
and the subtle color gradients of dyed textiles. The breakthrough came
not from fashion-specific work but from the ImageNet classification
challenge: Krizhevsky et al.\ \cite{krizhevsky2012imagenet} showed with
AlexNet that learned hierarchical features outperform any hand-designed
alternative at scale.

Subsequent architecture advances pushed accuracy further. Simonyan and
Zisserman \cite{simonyan2015vggnet} demonstrated that simply stacking
$3{\times}3$ convolutions to greater depth (VGGNet, 16--19 layers)
improves feature quality, while He et al.\ \cite{he2016resnet}
introduced residual shortcuts that allow networks of 50+ layers to train
without the vanishing-gradient problem. Both VGG16 and ResNet50 remain
popular feature extractors in retrieval literature because their
intermediate representations can be repurposed as fixed-length image
embeddings with minimal effort.

On the fashion side, the DeepFashion benchmark \cite{liu2016deepfashion}
established that CNN embeddings support large-scale cross-domain
retrieval across 50 clothing categories and 1000 attributes.
Metric-learning approaches using triplet loss
\cite{schroff2015facenet} can sharpen the embedding space further, but
they require carefully curated positive/negative pairs and
attribute-level labels that many real-world catalogs lack. For practical
deployment scenarios --- ours included --- frozen ImageNet features
provide a strong and zero-cost starting point that avoids the expense
of supervised fine-tuning \cite{yosinski2014transferable}.

Beyond accuracy, recommendation quality has multiple dimensions.
Ziegler et al.\ \cite{ziegler2005diversification} argue that users
prefer varied rather than repetitive recommendation lists, and propose
intra-list diversity (ILD) as a complementary metric. Ge et al.\
\cite{ge2010beyond} formalize catalog coverage --- the fraction of
catalog items that ever appear in a recommendation --- as a measure of
long-tail exposure. We adopt both metrics in our evaluation. For
indexing, we use the Ball Tree data structure
\cite{omohundro1989five}, which supports exact nearest-neighbor queries
on L2-normalized vectors and avoids the complexity of approximate
methods while remaining practical for catalogs of our scale ($\sim$44K
items).


% ─────────────────────────────────────────────────────────────────────────────
\section{Methodology}
\label{sec:method}

Fig.~\ref{fig:architecture} outlines the system pipeline, which has
four stages: preprocessing, feature extraction, similarity indexing,
and deployment through a web interface.

\begin{figure}[!t]
  \centering
  \includegraphics[width=\columnwidth]{system_architecture_paper.png}
  \caption{System pipeline. Raw product images are resized and
           normalized, passed through a frozen CNN backbone to produce
           feature vectors, L2-normalized, indexed with a Ball Tree KNN
           structure, and served through a Streamlit web application.}
  \label{fig:architecture}
\end{figure}

\subsection{Image Preprocessing}

Every catalog image is resized to $224 \times 224$ pixels using
bilinear interpolation and channel-normalized with the ImageNet
training-set statistics:
\begin{equation}
  \tilde{x}_c = \frac{x_c - \mu_c}{\sigma_c}, \quad c \in \{R, G, B\}
  \label{eq:normalize_pixel}
\end{equation}
where $\mu = [0.485,\,0.456,\,0.406]$ and
$\sigma = [0.229,\,0.224,\,0.225]$. No data augmentation is applied
because we are extracting fixed descriptors, not training any layers.

\subsection{Feature Extraction}

We evaluate three ImageNet-pretrained backbones, all kept frozen
(no gradient updates):

\textbf{ResNet18} \cite{he2016resnet}: The lightest backbone in our
comparison (11.7\,M parameters). We remove the final classification
layer and read the global average pooling (GAP) output, yielding a
512-dimensional embedding per image.

\textbf{VGG16} \cite{simonyan2015vggnet}: The heaviest backbone
(138\,M parameters). We tap the output of the first fully-connected
layer (FC1 + ReLU), producing a 4096-dimensional embedding. Many of
these dimensions encode spatial redundancies from the three FC layers,
which, as we show in Section~\ref{sec:results}, hurts retrieval
relative to the more compact ResNet representations.

\textbf{ResNet50} \cite{he2016resnet}: A mid-range backbone (25.6\,M
parameters) that reads the 2048-dimensional GAP output. Residual
connections allow it to build richer, more hierarchical feature
representations than ResNet18 while remaining far more parameter
efficient than VGG16.

All three models process images in batches of 64 on a single NVIDIA
T4 GPU (Google Colab). For the ResNet variants, the backbone output
has shape $(B, D, 1, 1)$ and is reshaped to $(B, D)$:
\begin{equation}
  \mathbf{f} = \text{reshape}\!\left(\text{GAP}(\mathbf{z}),\;[B, D]\right)
  \label{eq:reshape}
\end{equation}
where $D \in \{512, 2048\}$.

\subsection{L2 Normalization and Cosine-Equivalent Retrieval}

After extraction, every feature vector is projected onto the unit
hypersphere:
\begin{equation}
  \hat{\mathbf{f}} = \frac{\mathbf{f}}{\|\mathbf{f}\|_2}
  \label{eq:l2norm}
\end{equation}
This step is critical because it establishes an algebraic equivalence
between Euclidean distance and cosine similarity. For any two unit
vectors $\hat{\mathbf{f}}_a$ and $\hat{\mathbf{f}}_b$:
\begin{equation}
  \cos(\hat{\mathbf{f}}_a, \hat{\mathbf{f}}_b)
  = 1 - \frac{\|\hat{\mathbf{f}}_a - \hat{\mathbf{f}}_b\|_2^2}{2}
  \label{eq:cosine_equiv}
\end{equation}
Equation~\eqref{eq:cosine_equiv} lets us use a Ball Tree index
\cite{omohundro1989five} --- which natively supports Euclidean queries
--- to perform cosine-equivalent retrieval without materializing a full
$N \times N$ similarity matrix. We store all features as
\texttt{float32} arrays via \texttt{joblib} serialization, halving
memory relative to the default \texttt{float64}.

\subsection{Ball Tree KNN Index}

Given a query vector $\hat{\mathbf{q}}$, the top-$K$ nearest neighbors
are retrieved as:
\begin{equation}
  \mathcal{N}_K(\hat{\mathbf{q}})
  = \operatorname*{arg\,min}_{K}\;
    \|\hat{\mathbf{q}} - \hat{\mathbf{f}}_i\|_2, \quad i = 1,\ldots,N
  \label{eq:knn}
\end{equation}
The cosine similarity for each neighbor is then recovered from
Eq.~\eqref{eq:cosine_equiv}. We set $K{=}6$ throughout, following
standard practice in top-$K$ recommendation evaluation.


% ─────────────────────────────────────────────────────────────────────────────
\section{Dataset and Experimental Setup}
\label{sec:setup}

\subsection{Dataset}

The Myntra fashion dataset is a publicly available collection of Indian
e-commerce product images accompanied by structured metadata in a CSV
file (\texttt{styles.csv}). Each record contains fields for gender,
master category, sub-category, article type, base colour, season, year,
and product display name. Table~\ref{tab:dataset} summarizes the key
properties.

\begin{table}[!t]
\caption{Myntra Dataset Statistics}
\label{tab:dataset}
\centering
\begin{tabular}{lll}
\toprule
\textbf{Property}  & \textbf{Value}      & \textbf{Details} \\
\midrule
Total images       & 44,441              & All loaded successfully \\
Feature dim (R18)  & 512-D               & ResNet18 GAP output \\
Feature dim (VGG16)& 4096-D              & VGG16 FC1 output \\
Feature dim (R50)  & 2048-D              & ResNet50 GAP output \\
Master categories  & 7                   & Apparel, Accessories, etc. \\
Article types      & 142                 & Unique product types \\
Base colours       & 46                  & From \texttt{styles.csv} \\
Batch size         & 64                  & Feature extraction \\
Input resolution   & $224{\times}224$    & All backbones \\
\bottomrule
\end{tabular}
\end{table}

The seven master categories are heavily imbalanced: Apparel (21,392),
Accessories (11,274), Footwear (9,219), Personal Care (2,403), Free
Items (105), Sporting Goods (25), and Home (1). Apparel alone accounts
for 48\% of the catalog, which would dominate any uniformly random
evaluation sample. The color distribution is led by Black and Navy
Blue, and the gender split skews heavily toward Men's products.
Fig.~\ref{fig:cat_dist} visualizes the category breakdown.

\begin{figure}[!t]
  \centering
  \includegraphics[width=\columnwidth]{Master Category Distribution.png}
  \caption{Master category distribution in the Myntra dataset. The
           dominance of Apparel (48\%) motivates our stratified
           evaluation protocol.}
  \label{fig:cat_dist}
\end{figure}

\begin{figure}[!t]
  \centering
  \includegraphics[width=\columnwidth]{Confusion Matrix.png}
  \caption{Confusion matrix for ResNet50 on the stratified 166-query
           benchmark. Diagonal entries confirm that the majority of
           retrieved neighbors share the query's master category.}
  \label{fig:confusion}
\end{figure}

\subsection{Evaluation Protocol}

Sampling queries uniformly at random would allocate roughly half of
them to Apparel, making it impossible to assess whether the system
handles minority classes like Sporting Goods or Home. Instead, we use
stratified sampling: $\lfloor 200/7 \rfloor \approx 28$ queries are
drawn from each master category, yielding 166 balanced queries across
the seven classes.

For each query, the top-$K{=}6$ neighbors are retrieved (excluding the
query itself if it appears in the index). The predicted label is
determined by majority vote among the categories of the six returned
items. We compute four classification metrics against the ground-truth
labels from \texttt{styles.csv}:

\begin{itemize}
  \item \textbf{Accuracy}: fraction of queries where the majority-vote
        label matches the query's true category.
  \item \textbf{Precision@6}: fraction of retrieved items that share
        the query's category, averaged across all queries.
  \item \textbf{Recall@6}: fraction of same-category catalog items
        retrieved, averaged and capped at $K$.
  \item \textbf{F1-Score@6}: harmonic mean of Precision@6 and Recall@6.
\end{itemize}

We also measure retrieval coherence via Top-1 and average Top-5 cosine
similarity; intra-list diversity
$\text{DIV} = 1 - \text{ILS}$ where ILS is the mean pairwise cosine
similarity within each recommendation list
\cite{ziegler2005diversification}; and catalog coverage, defined as the
fraction of distinct catalog items that appear in any recommendation
across 500 random queries \cite{ge2010beyond}.


% ─────────────────────────────────────────────────────────────────────────────
\section{Results and Discussion}
\label{sec:results}

\subsection{Backbone Comparison}

Table~\ref{tab:comparison} reports all metrics for the three backbones
on the stratified 166-query benchmark.

\begin{table}[!t]
\caption{Backbone Comparison --- Myntra Dataset\\
(166 stratified queries, $K{=}6$, frozen ImageNet weights)}
\label{tab:comparison}
\centering
\renewcommand{\arraystretch}{1.2}
\begin{tabular}{lccccc}
\toprule
\textbf{Model} & \textbf{Acc} & \textbf{P@6} & \textbf{R@6} &
\textbf{F1@6} & \textbf{Top-1} \\
\midrule
ResNet18 (512-D)   & 0.759 & 0.762 & 0.762 & 0.762 & 0.924 \\
VGG16 (4096-D)     & 0.729 & 0.762 & 0.762 & 0.762 & 0.878 \\
\textbf{ResNet50 (2048-D)} &
\textbf{0.777} & \textbf{0.789} & \textbf{0.789} &
\textbf{0.789} & \textbf{0.940} \\
\midrule
\textbf{Model} & \multicolumn{2}{c}{\textbf{Avg Top-5}} &
\multicolumn{2}{c}{\textbf{Diversity}} & \textbf{Params} \\
\midrule
ResNet18   & \multicolumn{2}{c}{0.906} & \multicolumn{2}{c}{0.111} & 11.7M \\
VGG16      & \multicolumn{2}{c}{0.852} & \multicolumn{2}{c}{\textbf{0.163}} & 138M \\
\textbf{ResNet50} & \multicolumn{2}{c}{\textbf{0.925}} &
\multicolumn{2}{c}{0.088} & 25.6M \\
\bottomrule
\multicolumn{6}{l}{\small Bold = best in column. Catalog coverage = 6.4\% for all three models.}
\end{tabular}
\end{table}

\begin{figure}[!t]
  \centering
  \includegraphics[width=\columnwidth]{Precision Recall F1.png}
  \caption{Accuracy and F1-Score@6 across the three backbones. ResNet50
           leads on both metrics despite having only 25.6\,M parameters
           --- roughly a fifth of VGG16's 138\,M.}
  \label{fig:comparison}
\end{figure}

ResNet50 wins every classification and similarity column. Its Accuracy
of 0.777 and F1@6 of 0.789 are reached without any fashion-specific
training whatsoever, confirming that residual features learned on
ImageNet transfer effectively to fashion product images. The gain over
ResNet18 is modest (+1.8 pp accuracy, +2.7 pp F1), suggesting that
doubling the embedding dimension from 512 to 2048 provides a measurable
but not dramatic improvement.

VGG16 presents a more interesting case. Despite commanding 138\,M
parameters --- roughly 12 times ResNet18 and over five times
ResNet50 --- it trails both ResNets on accuracy and ties ResNet18 on
F1@6. We attribute this to the nature of VGG16's 4096-D
representation: the three fully-connected layers at the end of the
network encode substantial spatial redundancy that helps ImageNet
classification but actively hurts retrieval, where compact,
discriminative embeddings are preferred. ResNet50's 2048-D features,
shaped by residual connections that preserve gradient flow across 50
layers, capture category-level semantics more efficiently than VGG16
manages with six times the parameter budget.

Interestingly, VGG16 does produce the most diverse recommendation lists
(Diversity 0.163 vs.\ 0.088 for ResNet50). The lower feature precision
of VGG16 means that its nearest-neighbor sets span a wider visual range
--- the system recommends a broader variety of items rather than six near
duplicates. Whether this trade-off is desirable depends on the UX
context: a ``more like this'' widget benefits from tight similarity
(ResNet50), while an ``explore this style'' carousel might prefer the
variety that VGG16 provides.

All three models produce the same catalog coverage (6.4\%), indicating
that each model retrieves from a small, stable set of near-neighbors
and rarely surfaces long-tail items. Improving coverage through
query-time diversification or approximate indexing with randomized
tie-breaking is a direction for future investigation.

\subsection{Feature Space Analysis}

\begin{figure}[!t]
  \centering
  \includegraphics[width=\columnwidth]{Free Space-PCA Projection.png}
  \caption{PCA projection of ResNet50 embeddings (1,000 items, colored
           by master category). Footwear and Accessories form tight
           clusters; Apparel spans a wider region.}
  \label{fig:pca}
\end{figure}

\begin{figure}[!t]
  \centering
  \includegraphics[width=\columnwidth]{tSNE_Visualization.png}
  \caption{t-SNE projection of ResNet50 embeddings (2,000 items).
           Tighter, more separated clusters than PCA confirm strong
           local neighborhood structure organized by category.}
  \label{fig:tsne}
\end{figure}

PCA and t-SNE \cite{vandermaaten2008tsne} projections of the ResNet50
feature space (Figs.~\ref{fig:pca} and \ref{fig:tsne}) both reveal
cluster structure aligned with product categories. Footwear and
Accessories form compact, well-separated groups --- unsurprisingly,
since shoes and watches have distinctive shapes (sole profiles, dial
faces) that ImageNet features capture well. Apparel occupies a broader
region of the embedding space, reflecting the high visual diversity
within that category: tops, bottoms, dresses, and ethnic wear all
carry the same ``Apparel'' label yet look nothing alike. This
observation partly explains why Apparel queries are more likely to
produce cross-category errors in the confusion matrix
(Fig.~\ref{fig:confusion}).

\subsection{Similarity and Retrieval Quality}

\begin{figure}[!t]
  \centering
  \includegraphics[width=\columnwidth]{Top-K Similarity Distributions.png}
  \caption{Distribution of Top-1, Avg Top-3, and Avg Top-5 cosine
           similarity for ResNet50. The mass of the Top-1 distribution
           above 0.9 confirms that the nearest neighbor is almost
           always a strong visual match.}
  \label{fig:topk}
\end{figure}

ResNet50's Top-1 cosine similarity of 0.940 (Fig.~\ref{fig:topk})
means that, for the typical query, the closest item in the catalog
shares over 94\% of its visual signature. The average Top-5 similarity
of 0.925 shows only a gentle decline across the recommendation list,
confirming that all six returned items maintain strong visual coherence
with the query.

The intra-list diversity of 0.088 for ResNet50 tells the same story
from the opposite direction: the six recommendations are highly
similar to each other, differing by less than 9\% on average.
For a ``find similar items'' use case, this tight clustering is
exactly the desired behavior.


% ─────────────────────────────────────────────────────────────────────────────
\section{Web Application}
\label{sec:app}

The system is deployed as a Streamlit \cite{streamlit2024} application
with three navigation modes, selectable from a sidebar. The interface
follows a dark-theme aesthetic with gold accent colors, drawing on
visual conventions from luxury fashion branding. Custom CSS injected
via Streamlit's \texttt{unsafe\_allow\_html} parameter controls
typography (Cormorant Garamond for headings, Outfit for body text),
card layouts with animated gold borders for KPI metrics, and smooth
micro-animations on hover states. This level of visual polish
demonstrates that a research prototype can be wrapped in a
near-production-quality interface with relatively little additional
effort.

\textbf{Browse Catalog} (Fig.~\ref{fig:ui_browse}) presents a grid of
randomly sampled product images. Clicking any item triggers a KNN query
and displays the top-6 visual matches alongside similarity scores,
product names, and category labels. Similarity scores are color-coded
(green above 85\%, gold above 70\%, red below) to give an immediate
visual quality signal.

\textbf{Upload \& Discover} (Fig.~\ref{fig:ui_upload}) accepts a
user-uploaded image, runs it through the selected backbone's
preprocessing and forward pass in real time, and retrieves catalog
matches. This mode enables a ``shop from a photo'' workflow: a user
can photograph an item they like and instantly see similar products
from the Myntra catalog.

\textbf{Analytics Dashboard} provides three tabs: (i)~Dataset \& EDA
--- category, colour, gender, season, and year distributions, plus
PCA and t-SNE projections of the feature space; (ii)~Classification
Metrics --- accuracy, precision, recall, F1, per-category bar charts,
confusion matrix, and a full classification report; (iii)~Recommendation
Metrics --- Top-K similarity histograms, a diversity gauge, a catalog
coverage curve, and a system scorecard.

The application supports real-time model switching: the sidebar exposes
a radio button for ResNet18, VGG16, or ResNet50. Selecting a model
loads the corresponding precomputed feature matrix and KNN index from
a serialized \texttt{.pkl} file (generated offline via \texttt{joblib}).
This design avoids recomputing features at query time, keeping response
latency under one second for catalog browsing. For uploaded images, a
single forward pass through the selected backbone adds roughly two
seconds of latency on a CPU-only machine.

The full pipeline is built with PyTorch \cite{paszke2019pytorch} for
feature extraction and scikit-learn \cite{pedregosa2011sklearn} for KNN
indexing and metric computation. The application requires no GPU at
inference time, making it deployable on commodity hardware.

\begin{figure}[!t]
  \centering
  \includegraphics[width=\columnwidth]{Streamlit interface.png}
  \caption{Browse Catalog mode. The selected query item appears on the
           left; its six nearest visual matches with similarity scores
           are displayed on the right.}
  \label{fig:ui_browse}
\end{figure}

\begin{figure}[!t]
  \centering
  \includegraphics[width=\columnwidth]{Streamlit interface1.png}
  \caption{Upload \& Discover mode. A user-uploaded product image
           retrieves visually similar catalog items in real time.}
  \label{fig:ui_upload}
\end{figure}


% ─────────────────────────────────────────────────────────────────────────────
\section{Conclusion}
\label{sec:conclusion}

\subsection{Summary}

We compared three frozen ImageNet backbones --- ResNet18, VGG16, and
ResNet50 --- for content-based fashion retrieval on 44,441 Myntra
product images. Using a stratified evaluation protocol that samples
equally from seven product categories, we found that ResNet50 achieves
the best results across all classification and similarity metrics
(Accuracy 0.777, F1-Score@6 0.789, Top-1 Similarity 0.940) without
any domain-specific fine-tuning. VGG16, despite its 138\,M parameters,
underperforms both ResNets on retrieval accuracy but produces the
most diverse recommendation lists (Diversity 0.163). ResNet18 offers
a reasonable accuracy--efficiency trade-off for resource-constrained
deployments.

The accompanying Streamlit application integrates visual recommendation,
image-upload search, and a multi-tab analytics dashboard into a single
deployable system that runs without a GPU at inference time.

\subsection{Limitations}

Several limitations of our study deserve explicit mention. First, the
Myntra catalog has extreme category imbalance: Home contains exactly
one item and Sporting Goods contains only 25. Although our stratified
protocol assigns equal queries to each category, the Ball Tree index
still draws neighbors from these tiny pools, so minority-category
retrieval effectively reduces to near-exhaustive search rather than
efficient approximate lookup.

Second, all three backbones are frozen at ImageNet weights. While this
zero-shot approach works well for coarse category-level retrieval, it
cannot capture fashion-specific attributes such as neckline shape,
sleeve length, or fabric texture --- attributes that matter for
fine-grained ``more like this'' recommendations. Fine-tuning on a
fashion-attribute dataset would likely close this gap but requires
labeled training data that the Myntra CSV does not supply.

Third, the evaluation is based entirely on category-level ground truth.
A recommendation that returns a red cotton kurta when the query is a
blue silk kurta counts as correct if both are labeled Apparel, even
though a real shopper would consider them very different. Attribute-level
evaluation would give a more nuanced picture of retrieval quality.

Finally, the Ball Tree index does not scale gracefully beyond
$\sim$100K items. Production systems typically use approximate
nearest-neighbor libraries such as FAISS \cite{johnson2019faiss} or
ScaNN, which trade a small amount of recall for orders-of-magnitude
speedup. Our exact-search setup is appropriate for a 44,441-item
catalog but would need replacement for million-scale deployment.

\subsection{Future Work}

Three extensions follow naturally from the limitations above:
(i)~fine-tuning the backbone on fashion-specific attributes using
metric learning with triplet or contrastive loss;
(ii)~replacing the Ball Tree with an approximate nearest-neighbor
index (e.g., FAISS) to support catalogs at million-item scale;
(iii)~combining visual features with textual product descriptions
(product display name, article type, colour) for multimodal
recommendation.


% ─────────────────────────────────────────────────────────────────────────────
\begin{thebibliography}{17}

\bibitem{he2016resnet}
K.~He, X.~Zhang, S.~Ren, and J.~Sun,
``Deep residual learning for image recognition,''
in \emph{Proc. IEEE CVPR}, 2016, pp. 770--778.

\bibitem{simonyan2015vggnet}
K.~Simonyan and A.~Zisserman,
``Very deep convolutional networks for large-scale image recognition,''
in \emph{Proc. ICLR}, 2015.

\bibitem{krizhevsky2012imagenet}
A.~Krizhevsky, I.~Sutskever, and G.~E.~Hinton,
``ImageNet classification with deep convolutional neural networks,''
in \emph{Adv. Neural Inf. Process. Syst.}, vol.~25, 2012, pp. 1097--1105.

\bibitem{liu2016deepfashion}
Z.~Liu, P.~Luo, S.~Qiu, X.~Wang, and X.~Tang,
``DeepFashion: Powering robust clothes recognition and retrieval with
rich annotations,''
in \emph{Proc. IEEE CVPR}, 2016, pp. 1096--1104.

\bibitem{yosinski2014transferable}
J.~Yosinski, J.~Clune, Y.~Bengio, and H.~Lipson,
``How transferable are features in deep neural networks?''
in \emph{Adv. Neural Inf. Process. Syst.}, vol.~27, 2014, pp. 3320--3328.

\bibitem{mcauley2015image}
J.~McAuley, C.~Targett, Q.~Shi, and A.~van den Hengel,
``Image-based recommendations on styles and substitutes,''
in \emph{Proc. 38th ACM SIGIR}, 2015, pp. 43--52.

\bibitem{lops2011content}
P.~Lops, M.~de~Gemmis, and G.~Semeraro,
``Content-based recommender systems: State of the art and trends,''
in \emph{Recommender Systems Handbook}. Springer, 2011, pp. 73--105.

\bibitem{schroff2015facenet}
F.~Schroff, D.~Kalenichenko, and J.~Philbin,
``FaceNet: A unified embedding for face recognition and clustering,''
in \emph{Proc. IEEE CVPR}, 2015, pp. 815--823.

\bibitem{ziegler2005diversification}
C.-N.~Ziegler, S.~M.~McNee, J.~A.~Konstan, and G.~Lausen,
``Improving recommendation lists through topic diversification,''
in \emph{Proc. 14th WWW}, 2005, pp. 22--32.

\bibitem{ge2010beyond}
M.~Ge, C.~Delgado-Battenfeld, and D.~Jannach,
``Beyond accuracy: Evaluating recommender systems by coverage and
serendipity,''
in \emph{Proc. 4th ACM RecSys}, 2010, pp. 257--260.

\bibitem{lowe2004distinctive}
D.~G.~Lowe,
``Distinctive image features from scale-invariant keypoints,''
\emph{Int. J. Comput. Vis.}, vol. 60, no. 2, pp. 91--110, 2004.

\bibitem{omohundro1989five}
S.~M.~Omohundro,
``Five balltree construction algorithms,''
\emph{Int. Comput. Sci. Inst. Tech. Rep.}, 1989.

\bibitem{vandermaaten2008tsne}
L.~van der Maaten and G.~Hinton,
``Visualizing data using t-SNE,''
\emph{J. Mach. Learn. Res.}, vol.~9, pp. 2579--2605, 2008.

\bibitem{paszke2019pytorch}
A.~Paszke \emph{et al.},
``PyTorch: An imperative style, high-performance deep learning library,''
in \emph{Adv. Neural Inf. Process. Syst.}, vol.~32, 2019, pp. 8026--8037.

\bibitem{pedregosa2011sklearn}
F.~Pedregosa \emph{et al.},
``Scikit-learn: Machine learning in Python,''
\emph{J. Mach. Learn. Res.}, vol.~12, pp. 2825--2830, 2011.

\bibitem{johnson2019faiss}
J.~Johnson, M.~Douze, and H.~J\'{e}gou,
``Billion-scale similarity search with GPUs,''
\emph{IEEE Trans. Big Data}, vol.~7, no.~3, pp. 535--547, 2021.

\bibitem{streamlit2024}
Streamlit Inc.,
``Streamlit --- A faster way to build and share data apps,''
\url{https://streamlit.io}, 2024.

\end{thebibliography}

\end{document}
